% Insert after the ordinary and odd fourth-moment entrances.
% This is the q-anchored/W_2 interface, not the ordinary no-pair lemma.
\subsection{An anchored packet and its reversed unit row}
\label{sec:odd-common:q-packet}

\begin{lemma}[A correlated first row and the \(W_2\) cancellation]
\label{lem:odd-common:q-packet}
Let \(F\) be an actual terminal, \(R=r(B)u\), and suppose its actual
first row has been proved in the form
\begin{equation}\label{eq:odd-common:q-first-row}
 TF\le e^{\alpha}(\bar f_2+BR+D),\qquad D\ge0.
\end{equation}
For a terminal \(F\le q e^{R-\Gamma}\), \(\alpha=-\Gamma\).
A square row with prefactor \(m^2e^{-\Gamma}\) instead has
\(\alpha=2\log m-\Gamma\); that scalar is kept throughout.
Set \(D_i=\E B^iR\) and \(\gamma_j=\beta_j-D_{j-2}\).
Finite propagation and the positive path convolution give
\begin{equation}\label{eq:odd-common:q-finite}
 T^LF\le e^{\alpha}\left[
 1+G_{L+1}(B)u+B^LR+D\bar f_{L-1}
       -\sum_{j=3}^{L+1}\gamma_j\bar f_{L+1-j}\right].
\end{equation}
For \(L\ge2\), if the displayed \(\gamma_j\) are nonnegative,
\begin{equation}\label{eq:odd-common:q-long}
\begin{gathered}
 T^LF\le e^{\alpha}[1+u+c_L+BQ_L(B)u],\qquad
 c_L=2D-\gamma_{L+1},\quad Q_L=G_L+z^{L-1}r,\\
 \E BQ_L(B)u=-\gamma_{L+2}.
\end{gathered}
\end{equation}
The mean identity requires the same actual moment at index \(L+2\).

Under complete terminal reversal, \(\bar R=-r(-B)u\) and
\begin{equation}\label{eq:odd-common:q-reversal}
 \E\bar R=\E R,\quad\|\bar R\|=\|R\|,
 \quad\E(u\bar R)=-\E(uR).
\end{equation}
If the first-row proof used
\(D\ge\max\{\E R+\E(uR)+\sqrt{d\E R^2}+c_T\E R^2,0\}\),
the same proof and the same \(D\) give
\(T\bar F\le e^{\alpha}[Tp+B\bar R+D]\).
The coefficient \(c_T\) is the one proved for the actual terminal.

Suppose the chosen root pair has the full envelope
\(P_r\le e^{-\Gamma_r+BV_r}\), including its actual constant.
The resulting complete certificates are
\begin{equation}\label{eq:odd-common:q-cancel}
\begin{array}{c|c|c}
 \text{actual message}&a&\text{complete }V\\ \hline
 f_2P_rT\bar F&\alpha+D-\Gamma_r&\bar R+V_r\\
 W_2P_rT^LF,\ L\ge2&\alpha+c_L-\Gamma_r&(Q_L-1)u+V_r.
\end{array}
\end{equation}
They mean \(\log H\le a+BV\), with fourth exponent
\(4(a+\E BV)+(17/4)\|V\|^2\).
For the second row its scalar-plus-mean, before the actual root
mean \(\E BV_r\), is
\(\alpha-\Gamma_r+2D+d-\gamma_{L+1}-\gamma_{L+2}\).
\end{lemma}

\begin{proof}
Apply \(T^{L-1}\) to \eqref{eq:odd-common:q-first-row} and use
the bare finite Duhamel identity
\eqref{eq:oc:finite-duhamel}. Substituting
\eqref{eq:Q188:conv} gives
\eqref{eq:odd-common:q-finite}. Retaining the final
\(\gamma_{L+1}\bar f_0\) and using
\(\bar f_{L-1}\le2\) gives \eqref{eq:odd-common:q-long}; the
two spectral means give its next-index identity and
\eqref{eq:odd-common:q-reversal}.
In the first-row calculation
\(T(qR)=\E R+\E(uR)+BR+B(uR)\).
Reversal changes its mean sum to
\(\E\bar R-\E(u\bar R)\), which is the same number;
the row-product bound remains \(\sqrt{d\E R^2}\).
Finally \(f_2=1-u+Bu\) cancels both \(u\) and \(Bu\) in
\(Tp=1+u-Bu\), whereas \(W_2=1-u-Bu\) cancels \(u\) and
subtracts one \(u\) from the primitive in the long row.
Apply \(\log y\le y-1\) to these positive factors and their
positive majorants. This gives the displayed full vectors,
without a new unit-arm curvature reserve.
\end{proof}
